\section{Introduction}

Time-series foundation models promise forecasts without training: a pre-trained network receives the history of a series and
returns quantile forecasts for any horizon. TimesFM~3, released by Google Research on 31 August 2026 \citep{jain2026timesfm3}, extends
the TimesFM line \citep{das2024timesfm} to multivariate inputs and covariates and is reported to lead the public benchmarks
GIFT-Eval and fev-bench \citep{aksu2024gifteval,shchur2025fevbench}. For a practitioner the questions are concrete: will the model
beat the classical methods already in use on \emph{my} series, by how much, where will it fail, and is the extra compute worth it?

Public benchmarks answer these questions only in part. Their datasets are years old and widely redistributed, and the training
corpora of foundation models are large, partly undisclosed and drawn from the same public sources; a recent survey finds that only a
small fraction of the datasets used to evaluate these models never appeared in any pre-training corpus \citep{meyer2025leakage}. A
benchmark advantage may therefore mix forecasting skill with exposure. The reported comparisons also depend on inference options
whose role is rarely visible in a leaderboard: symmetric averaging, clipping at zero, or, as we show, letting a competitor share
information across up to 100 series at once.

This paper evaluates TimesFM~3 zero-shot under a protocol frozen before the post-release test data were downloaded, with data that
\emph{post-date the model's release}. We ask seven questions: whether the published results reproduce (RQ1); whether the advantage
holds on post-release data and leaves a trace of memorisation (RQ2); where zero-shot forecasting fails under controlled conditions
(RQ3); whether the quantiles are calibrated (RQ4); when multivariate inputs help (RQ5); how much the fixed inference steps
contribute (RQ6); and whether the model is worth its compute against classical methods (RQ7).

\paragraph{Contributions.}
\begin{itemize}
\item A \emph{leakage-aware design}: real data published after the release (six domains, 104 series, 4{,}068 series-windows), a
Wikipedia probe that compares the same pages inside and outside the model's documented pageview corpus, a prospective check with
forecasts sealed by hash before their outcomes existed, and controlled synthetic data, all scored with paired contrasts, a two-way
cluster bootstrap and Holm corrections within pre-declared families.
\item A \emph{reproduction} of Google's per-task fev-bench results and an explanation of why our Chronos-2 results differ from the
published ones (cross-learning), which shows that the published comparison gives Chronos-2 an option TimesFM~3 does not have.
\item A \emph{reading of the inference code} of TimesFM~3, which has no technical report, with a configuration audit and ablations of
its fixed steps: linear detrending behind an $R^2$ gate, iterative normalisation of masked horizon tokens, stitching, and the
inference options.
\item Evidence on \emph{when not to use} the model: near-random-walk series, saturating trends, and recent regime changes.
\end{itemize}

\begin{figure}[!t]
\centering
\pgfplotstableread{
x y em ep
1.0924 0.14 0.0653 0.0738
1.1022 1.14 0.0838 0.0925
0.7761 2.14 0.1402 0.1455
0.8082 3.14 0.0461 0.0406
0.8722 4.14 0.0941 0.0779
0.8919 5.14 0.0674 0.0563
}\domensmase
\pgfplotstableread{
x y em ep
1.0909 -0.14 0.0486 0.0738
1.0481 0.86 0.0832 0.1006
0.4933 1.86 0.1325 0.1685
0.7743 2.86 0.0366 0.0321
0.8006 3.86 0.0626 0.0594
0.8154 4.86 0.0574 0.0546
}\domenssql
\pgfplotstableread{
x y em ep
1.0130 0.14 0.0520 0.0536
1.0657 1.14 0.0782 0.0980
0.9167 2.14 0.0868 0.0792
0.9563 3.14 0.0279 0.0278
0.9794 4.14 0.0491 0.0487
0.9600 5.14 0.0481 0.0475
}\domchrmase
\pgfplotstableread{
x y em ep
1.0082 -0.14 0.0436 0.0459
1.0463 0.86 0.0730 0.0980
0.9248 1.86 0.0586 0.0592
0.9581 2.86 0.0246 0.0249
0.9810 3.86 0.0449 0.0430
0.9614 4.86 0.0439 0.0429
}\domchrsql
\pgfplotstableread{
x y em ep
1.0859 0.14 0.0675 0.0808
1.1339 1.14 0.1020 0.1363
0.8329 2.14 0.1782 0.1771
0.5846 3.14 0.0860 0.0776
0.6098 4.14 0.0808 0.0862
0.1835 5.14 0.0482 0.0713
}\domrwmase
\pgfplotstableread{
x y em ep
1.0912 -0.14 0.0619 0.0699
1.0634 0.86 0.1085 0.1504
0.4300 1.86 0.1246 0.1679
0.5162 2.86 0.0713 0.0669
0.4466 3.86 0.0527 0.0628
0.1692 4.86 0.0416 0.0587
}\domrwsql
\pgfplotsset{domain ticks/.style={ytick={0,1,2,3,4,5}, yticklabels={{Yields and oil},{Exchange rates},{River flow},{Weather},{Electricity price},{Electricity load}}, ymin=-0.6, ymax=5.6}}

\begin{tikzpicture}
\begin{groupplot}[
  group style={group size=3 by 1, horizontal sep=0.45cm, yticklabels at=edge left},
  vz axis, width=0.36\textwidth, height=6.4cm, xmode=log, log basis x=10, domain ticks,
  yticklabel style={font=\footnotesize, text=vzink, align=right},
  xticklabel style={/pgf/number format/fixed}, clip=false]
\nextgroupplot[title={vs classical ensemble}, xmin=0.33, xmax=1.4,
  xtick={0.4,0.6,0.8,1}, xticklabels={0.4,0.6,0.8,1}]
\fill[vzbluewash] (axis cs:0.33,-0.6) rectangle (axis cs:1,5.6);
\draw[vzaxis, line width=0.4pt] (axis cs:0.33,1.5) -- (axis cs:1.4,1.5);
\draw[vzink2, line width=0.6pt] (axis cs:1,-0.6) -- (axis cs:1,5.6);
\addplot[vz dots, mark=*, color=vzblue] table[x=x, y=y, x error minus=em, x error plus=ep] {\domensmase};
\addplot[vz dots, mark=diamond*, mark size=2.6pt, color=vzorange] table[x=x, y=y, x error minus=em, x error plus=ep] {\domenssql};
\node[font=\scriptsize, text=vzblue!70!black, anchor=north west, align=left] at (axis cs:0.345,5.55) {TimesFM~3\\better};
\nextgroupplot[title={vs Chronos-2}, xmin=0.8, xmax=1.25,
  xtick={0.8,0.9,1,1.1,1.2}, xticklabels={0.8,0.9,1,1.1,1.2}]
\fill[vzbluewash] (axis cs:0.8,-0.6) rectangle (axis cs:1,5.6);
\draw[vzaxis, line width=0.4pt] (axis cs:0.8,1.5) -- (axis cs:1.25,1.5);
\draw[vzink2, line width=0.6pt] (axis cs:1,-0.6) -- (axis cs:1,5.6);
\addplot[vz dots, mark=*, color=vzblue] table[x=x, y=y, x error minus=em, x error plus=ep] {\domchrmase};
\addplot[vz dots, mark=diamond*, mark size=2.6pt, color=vzorange] table[x=x, y=y, x error minus=em, x error plus=ep] {\domchrsql};
\nextgroupplot[title={vs random walk}, xmin=0.11, xmax=1.4,
  xtick={0.125,0.25,0.5,1}, xticklabels={0.125,0.25,0.5,1}]
\fill[vzbluewash] (axis cs:0.11,-0.6) rectangle (axis cs:1,5.6);
\draw[vzaxis, line width=0.4pt] (axis cs:0.11,1.5) -- (axis cs:1.4,1.5);
\draw[vzink2, line width=0.6pt] (axis cs:1,-0.6) -- (axis cs:1,5.6);
\addplot[vz dots, mark=*, color=vzblue] table[x=x, y=y, x error minus=em, x error plus=ep] {\domrwmase};
\addplot[vz dots, mark=diamond*, mark size=2.6pt, color=vzorange] table[x=x, y=y, x error minus=em, x error plus=ep] {\domrwsql};
\end{groupplot}
\node[font=\footnotesize, text=vzink2, anchor=north] at ($(group c2r1.south)+(0,-0.55cm)$)
  {error of TimesFM~3 / error of the reference (log scale; left of 1: TimesFM~3 better)};
\node[font=\footnotesize, text=vzink2, anchor=north] at ($(group c2r1.south)+(0,-1.0cm)$)
  {\tikz[baseline=-0.6ex]{\draw[vzblue, line width=1pt] (-0.22,0) -- (0.22,0); \fill[vzblue] (0,0) circle (2.1pt);}~MASE (point forecast)\qquad
   \tikz[baseline=-0.6ex]{\draw[vzorange, line width=1pt] (-0.22,0) -- (0.22,0); \pgfsetfillcolor{vzorange}\pgfsetstrokecolor{vzorange}\pgfsetplotmarksize{2.6pt}\pgfuseplotmark{diamond*}}~SQL (probabilistic forecast)\qquad
   bars: 95\% interval};
\end{tikzpicture}
\caption{TimesFM~3 on data published after its release, domain by domain (Tier B, univariate mode, 2026-09-01 to 10-20). Each
point is the ratio of TimesFM~3's error to that of a reference on the same series and windows; bars are two-way cluster
bootstrap 95\% intervals within the domain (10--20 series clusters, cities for weather, and 7--8 weeks per domain, so intervals are wide). It wins clearly
against the classical ensemble on the four structured domains, narrowly against Chronos-2, and has a higher MASE than the random
walk on exchange rates, Treasury yields and oil prices (below the grey line), as do the point estimates of the other foundation models. On river flow the
pre-registered ensemble is poorly specified; against a re-specified one the SQL ratio is 0.742 instead of 0.493
(Section~\ref{sec:robustness}). A descriptive breakdown, not
pre-registered; the confirmatory tests pool the domains (Figure~\ref{fig:forest}, Table~\ref{tab:familyp}).}
\label{fig:domains}
\end{figure}


\paragraph{Main findings.} On a 30-task sample the published results largely reproduce, including the lead over Chronos-2. On the four structured post-release domains TimesFM~3 beats a classical ensemble
clearly (\PoneMase{} in MASE, \PoneSql{} in SQL) and Chronos-2 narrowly (about 5\%); against TimesFM~2.5 the SQL difference is
resolved and the MASE difference is not. The margin over classical methods is, descriptively, slightly smaller than on the benchmark, and the Wikipedia
probe neither resolves a difference nor identifies an effect of memorisation. In two overlapping, sealed prospective weeks the
quantile-loss advantage over the ensemble is confirmed, but on load, prices and river flow the point-forecast advantage is not
(driven mainly by electricity prices, descriptively), and no difference from Chronos-2 or TimesFM~2.5 is resolved. On exchange rates, interest rates and oil prices it is worse than the random walk, and the point estimates of the other foundation models are also above parity (Figure~\ref{fig:domains}), and in
synthetic data it fails on a saturating process (about four times the error of the best classical model) and after recent upward
level shifts (in an exploratory grid of 16 settings of one process, worse in 12, and by 17\% to 194\% in the nine with a shift of two
or more standard deviations at most 96 steps before the origin). In synthetic data,
multivariate mode helps when there is a strong cross-series signal and is equivalent within $\pm$3\% in the two controls without
one. The fixed inference steps matter little on
average.

The rest of the paper reviews related work (Section~\ref{sec:related}), explains how TimesFM~3 works (Section~\ref{sec:mechanism}),
describes the protocol and data (Sections~\ref{sec:protocol}--\ref{sec:data}), reports the results
(Sections~\ref{sec:reproduction}--\ref{sec:prospective}) and discusses their practical meaning and limits
(Sections~\ref{sec:discussion}--\ref{sec:limitations}).
